
\documentclass{article}
\usepackage{colortbl}
\usepackage{makecell}
\usepackage{multirow}
\usepackage{supertabular}

\begin{document}

\newcounter{utterance}

\centering \large Interaction Transcript for game `cladder', experiment `full\_v1.5\_default', episode 8830 with Qwen3.5{-}4B.
\vspace{24pt}

{ \footnotesize  \setcounter{utterance}{1}
\setlength{\tabcolsep}{0pt}
\begin{supertabular}{c@{$\;$}|p{.15\linewidth}@{}p{.15\linewidth}p{.15\linewidth}p{.15\linewidth}p{.15\linewidth}p{.15\linewidth}}
    \# & \multicolumn{2}{c}{Player} && \multicolumn{2}{c}{Game Master} \\
    \hline

    \theutterance \stepcounter{utterance}  
    & & & \multicolumn{4}{p{0.6\linewidth}}{
        \cellcolor[rgb]{0.9,0.9,0.9}{
            \makecell[{{p{\linewidth}}}]{
                \texttt{\tiny{[P1$\langle$GM]}}
                \texttt{You are an expert in causal inference. The following question is not a typical commonsense query, but rather a meticulously designed question created by a professor specializing in causal inference, intended to assess the students' mastery of the course content.} \\
\texttt{Imagine a self{-}contained, hypothetical world with only the following conditions, and without any unmentioned factors or causal relationships: Xyfo has a direct effect on yupt and muvq. Zupj has a direct effect on yupt. Yupt has a direct effect on muvq. Xyfo is unobserved. Method 1: We look at how yupt correlates with muvq case by case according to xyfo. Method 2: We look directly at how yupt correlates with muvq in general. To understand how yupt affects muvq, is it more correct to use the Method 1 than Method 2?} \\
\texttt{Start your answer with 'Yes' or 'No', followed by additional reasoning or evidence to support your explanation.} \\
            }
        }
    }
     \\ \\

    \theutterance \stepcounter{utterance}  
    & \multicolumn{4}{p{0.6\linewidth}}{
        \cellcolor[rgb]{0.561,0.737,0.561}{
            \makecell[{{p{\linewidth}}}]{
                \texttt{\tiny{[P1$\rangle$GM]}}
                \texttt{Yes} \\
\\ 
\texttt{**Reasoning:**} \\
\\ 
\texttt{In this hypothetical world, the causal structure is defined as follows:} \\
\texttt{1.  \$X\_\{yfo\} \textbackslash{}rightarrow Y\_\{upt\}\$} \\
\texttt{2.  \$Z\_\{upj\} \textbackslash{}rightarrow Y\_\{upt\}\$} \\
\texttt{3.  \$Y\_\{upt\} \textbackslash{}rightarrow M\_\{uvq\}\$} \\
\\ 
\texttt{This forms a **collider** structure at \$Y\_\{upt\}\$ (specifically, \$X\_\{yfo\} \textbackslash{}rightarrow Y\_\{upt\} \textbackslash{}leftarrow Z\_\{upj\}\$), but more importantly, \$Y\_\{upt\}\$ acts as a **mediator** between the unobserved cause \$X\_\{yfo\}\$ and the outcome \$M\_\{uvq\}\$.} \\
\\ 
\texttt{*   **Analysis of Method 2 (General Correlation):** If we simply look at the correlation between \$Y\_\{upt\}\$ and \$M\_\{uvq\}\$ in general (without conditioning on \$X\_\{yfo\}\$), we observe the total effect of \$Y\_\{upt\}\$ on \$M\_\{uvq\}\$. However, because \$X\_\{yfo\}\$ is unobserved, the observed correlation between \$Y\_\{upt\}\$ and \$M\_\{uvq\}\$ is confounded by the fact that \$Y\_\{upt\}\$ is also caused by \$X\_\{yfo\}\$. While \$Y\_\{upt\}\$ is a direct cause of \$M\_\{uvq\}\$, the "general" correlation mixes the effect of \$Y\_\{upt\}\$ on \$M\_\{uvq\}\$ with the indirect pathway \$X\_\{yfo\} \textbackslash{}rightarrow Y\_\{upt\} \textbackslash{}rightarrow M\_\{uvq\}\$. Without knowing the specific value of \$X\_\{yfo\}\$, we cannot isolate the *direct* causal effect of \$Y\_\{upt\}\$ on \$M\_\{uvq\}\$ from the background noise introduced by the unobserved \$X\_\{yfo\}\$.} \\
\\ 
\texttt{*   **Analysis of Method 1 (Stratification/Conditioning):** By looking at the correlation between \$Y\_\{upt\}\$ and \$M\_\{uvq\}\$ "case by case according to \$X\_\{yfo\}\$," we are effectively conditioning on the unobserved confounder (or rather, the common cause of the mediator).} \\
\texttt{    *   If we could observe \$X\_\{yfo\}\$, we could calculate the conditional correlation \$P(M\_\{uvq\} | Y\_\{upt\}, X\_\{yfo\})\$.} \\
\texttt{    *   In this conditional world, the influence of \$X\_\{yfo\}\$ is held constant. Any remaining correlation between \$Y\_\{upt\}\$ and \$M\_\{uvq\}\$ represents the **direct causal effect** of \$Y\_\{upt\}\$ on \$M\_\{uvq\}\$, free from the confounding influence of \$X\_\{yfo\}\$.} \\
\\ 
\texttt{**Conclusion:**} \\
\texttt{To understand *how* \$Y\_\{upt\}\$ affects \$M\_\{uvq\}\$ (i.e., to estimate the direct causal effect \$\textbackslash{}beta\_\{Y \textbackslash{}rightarrow} \\
            }
        }
    }
    & & \\ \\

    \theutterance \stepcounter{utterance}  
    & & & \multicolumn{2}{p{0.3\linewidth}}{
        \cellcolor[rgb]{0.9,0.9,0.9}{
            \makecell[{{p{\linewidth}}}]{
                \texttt{\tiny{[GM$|$GM]}}
                \texttt{Yes} \\
            }
        }
    }
    & & \\ \\

    \theutterance \stepcounter{utterance}  
    & & & \multicolumn{2}{p{0.3\linewidth}}{
        \cellcolor[rgb]{0.9,0.9,0.9}{
            \makecell[{{p{\linewidth}}}]{
                \texttt{\tiny{[GM$|$GM]}}
                \texttt{yes} \\
            }
        }
    }
    & & \\ \\

    \theutterance \stepcounter{utterance}  
    & & & \multicolumn{2}{p{0.3\linewidth}}{
        \cellcolor[rgb]{0.9,0.9,0.9}{
            \makecell[{{p{\linewidth}}}]{
                \texttt{\tiny{[GM$|$GM]}}
                \texttt{game\_result = WIN} \\
            }
        }
    }
    & & \\ \\

\end{supertabular}
}

\end{document}
